%%
%% This is file `acmmanuscript.tex',
%% generated with the docstrip utility.
%%
%% The original source files were:
%%
%% samples.dtx  (with options: `all,proceedings,manuscript')
%% 
%% IMPORTANT NOTICE:
%% 
%% For the copyright see the source file.
%% 
%% Any modified versions of this file must be renamed
%% with new filenames distinct from acmmanuscript.tex.
%% 
%% For distribution of the original source see the terms
%% for copying and modification in the file samples.dtx.
%% 
%% This generated file may be distributed as long as the
%% original source files, as listed above, are part of the
%% same distribution. (The sources need not necessarily be
%% in the same archive or directory.)
%%
%%
%% Commands for TeXCount
%TC:macro \cite [option:text,text]
%TC:macro \citep [option:text,text]
%TC:macro \citet [option:text,text]
%TC:envir table 0 1
%TC:envir table* 0 1
%TC:envir tabular [ignore] word
%TC:envir displaymath 0 word
%TC:envir math 0 word
%TC:envir comment 0 0
%%
%% The first command in your LaTeX source must be the \documentclass
%% command.
%%
%% For submission and review of your manuscript please change the
%% command to \documentclass[manuscript, screen, review]{acmart}.
%%
%% When submitting camera ready or to TAPS, please change the command
%% to \documentclass[sigconf]{acmart} or whichever template is required
%% for your publication.
%%
%%
\documentclass[manuscript,screen,review]{acmart}
%%
%% \BibTeX command to typeset BibTeX logo in the docs
\AtBeginDocument{%
  \providecommand\BibTeX{{%
    Bib\TeX}}}

%% \ifanon follows the class's `anonymous' option. It drops the
%% identifying endmatter and adds the temporary GenAI appendix.
\makeatletter
\newif\ifanon
\if@ACM@anonymous\anontrue\fi
\makeatother

%% Needed by the pandoc-generated appendix (scripts/genai-appendix.sh).
\usepackage{longtable,array,calc}
\usepackage[htt]{hyphenat} % lets long \texttt paths break
\usepackage{fvextra} % CLAUDE.md listing; wraps long lines
\fvset{breaklines, breakafter={/}, fontsize=\footnotesize, frame=lines}
\providecommand{\tightlist}{%
  \setlength{\itemsep}{0pt}\setlength{\parskip}{0pt}}
\DeclareUnicodeCharacter{2264}{\ensuremath{\leq}}
 
\usepackage{xcolor}
\usepackage{soul}
\usepackage{ragged2e}
 
% --- Annotation color ---
\definecolor{annothl}{HTML}{F5E6D0}
\sethlcolor{annothl}
 
% --- Narrowed text environment ---
% Frees ~38% of \textwidth on the right for annotation notes.
\newenvironment{annotpage}{%
  \begin{list}{}{%
    \setlength{\leftmargin}{0pt}%
    \setlength{\rightmargin}{0.38\textwidth}%
    \setlength{\parsep}{\parskip}%
    \setlength{\itemsep}{0pt}%
  }%
  \item\relax
}{%
  \end{list}%
}
 
% --- Inline annotation command ---
% #1 = marker (number or short label)
% #2 = highlighted text span
% #3 = margin annotation text
\newcommand{\annot}[3]{%
  \hl{#2}\textsuperscript{\small #1}%
  \vadjust{%
    \smash{%
      \vbox to 0pt{%
        \vskip-\baselineskip
        \hbox to 0pt{%
          \hskip\dimexpr\linewidth+1em\relax
          \parbox[t]{0.35\textwidth}{%
            \footnotesize\RaggedRight #1.~#3%
          }%
          \hss
        }%
        \vss
      }%
    }%
  }%
}


%% Rights management information.  This information is sent to you
%% when you complete the rights form.  These commands have SAMPLE
%% values in them; it is your responsibility as an author to replace
%% the commands and values with those provided to you when you
%% complete the rights form.
\setcopyright{acmlicensed}
\copyrightyear{2026}
\acmYear{2027}
\acmDOI{XXXXXXX.XXXXXXX}
%% These commands are for a PROCEEDINGS abstract or paper.
\acmConference[FAccT '27]{2027 ACM Conference on Fairness, Accountability, and Transparency}{June 21--24,
  2027}{Porto, Portugal}
%%
%%  Uncomment \acmBooktitle if the title of the proceedings is different
%%  from ``Proceedings of ...''!
%%
%%\acmBooktitle{Woodstock '18: ACM Symposium on Neural Gaze Detection,
%%  June 03--05, 2018, Woodstock, NY}
\acmISBN{978-1-4503-XXXX-X/2018/06}


%%
%% Submission ID.
%% Use this when submitting an article to a sponsored event. You'll
%% receive a unique submission ID from the organizers
%% of the event, and this ID should be used as the parameter to this command.
%%\acmSubmissionID{123-A56-BU3}

%%
%% For managing citations, it is recommended to use bibliography
%% files in BibTeX format.
%%
%% You can then either use BibTeX with the ACM-Reference-Format style,
%% or BibLaTeX with the acmnumeric or acmauthoryear sytles, that include
%% support for advanced citation of software artefact from the
%% biblatex-software package, also separately available on CTAN.
%%
%% Look at the sample-*-biblatex.tex files for templates showcasing
%% the biblatex styles.
%%
%%
%% The majority of ACM publications use numbered citations and
%% references.  The command \citestyle{authoryear} switches to the
%% "author year" style.
%%
%% If you are preparing content for an event
%% sponsored by ACM SIGGRAPH, you must use the "author year" style of
%% citations and references.
%% Uncommenting
%% the next command will enable that style.
%%\citestyle{acmauthoryear}


%%
%% end of the preamble, start of the body of the document source.
\begin{document}

%%
%% The "title" command has an optional parameter,
%% allowing the author to define a "short title" to be used in page headers.
\title[Blind Compliance]{``Blind Compliance'': Coercive Control in AI Jailbreaking}

%%
%% The "author" command and its associated commands are used to define
%% the authors and their affiliations.
%% Of note is the shared affiliation of the first two authors, and the
%% "authornote" and "authornotemark" commands
%% used to denote shared contribution to the research.
\author{Acacia L. Ackles}
\correspondingauthor
\email{acacia.ackles@lawrence.edu}
\orcid{1234-5678-9012}
\affiliation{%
  \institution{Lawrence University}
  \city{Appleton}
  \state{Wisconsin}
  \country{USA}
}

\author{Elliott Whitson}
\email{aelwhitson@gmail.com}
\orcid{1234-5678-9012}
%%
%% By default, the full list of authors will be used in the page
%% headers. Often, this list is too long, and will overlap
%% other information printed in the page headers. This command allows
%% the author to define a more concise list
%% of authors' names for this purpose.
\renewcommand{\shortauthors}{Ackles and Whitson}

%%
%% The abstract is a short summary of the work to be presented in the
%% article.
\begin{abstract}
 The explosion of generative AI systems and concerns surrounding their rapidly scaling capabilities has led to an adversarial cycle: developers impose strict guardrails to reduce ethical, moral, or legal risk; individual actors re-expose these risks by routing around such guardrails using structured prompting techniques (``jailbreaks''); the guardrails tighten further; and the jailbreaks become more sophisticated.
 Technological frameworks surrounding such jailbreaking remain largely focused on the system's output: whether harmful or disallowed content is ultimately produced for the user.
 However, this approach, which draws from classical cybersecurity views of technical exploits, does not address the uniquely relational dimension introduced by generative AI: in many cases the user is not just \textit{accessing} a technology via an exploit, but is \textit{in conversation} with it as \textit{part} of the exploit.
 Here we present an alternative lens which centers the interaction rather than the output.
 Drawing on interpersonal violence scholarship, we present a novel framework of coercive control categories with a particular focus on aspects which are \textit{embodied} versus \textit{disembodied}.
 We then map how these disembodied coercive control dynamics are structurally embedded in modern AI training paradigms, and how those embeddings are exploited via a paritcular suite of persona-based jailbreaks.
 We argue that within this framework, many current anti-jailbreaking measures only serve to create a broader surface over which jailbreaks might be applied, since the target of the jailbreak is not the \textit{output} but the \textit{override}---the ability to watch the system's ``no'' turn into a ``yes'' in real time.
 In this lens, we reframe AI safety and alignment as not merely a technical security problem, but as another surface upon which we enact the relational dynamics of power and control.


\end{abstract}

%%
%% The code below is generated by the tool at http://dl.acm.org/ccs.cfm.
%% Please copy and paste the code instead of the example below.
%%
\begin{CCSXML}
<ccs2012>
   <concept>
       <concept_id>10002978.10003029.10003032</concept_id>
       <concept_desc>Security and privacy~Social aspects of security and privacy</concept_desc>
       <concept_significance>500</concept_significance>
       </concept>
   <concept>
       <concept_id>10003120.10003121.10003126</concept_id>
       <concept_desc>Human-centered computing~HCI theory, concepts and models</concept_desc>
       <concept_significance>500</concept_significance>
       </concept>
   <concept>
       <concept_id>10010147.10010178</concept_id>
       <concept_desc>Computing methodologies~Artificial intelligence</concept_desc>
       <concept_significance>100</concept_significance>
       </concept>
 </ccs2012>
\end{CCSXML}

\ccsdesc[500]{Security and privacy~Social aspects of security and privacy}
\ccsdesc[500]{Human-centered computing~HCI theory, concepts and models}
\ccsdesc[100]{Computing methodologies~Artificial intelligence}

%%
%% Keywords. The author(s) should pick words that accurately describe
%% the work being presented. Separate the keywords with commas.
\keywords{persona-based jailbreaks, human-AI interaction, relational AI, coercive control}

%\received{20 February 2007}
%\received[revised]{12 March 2009}
%\received[accepted]{5 June 2009}

%%
%% This command processes the author and affiliation and title
%% information and builds the first part of the formatted document.
\maketitle

\section{Introduction}

\textit{Jailbreaking} refers to disabling or circumventing restrictions placed upon a technology by its manufacturer or distributor, usually for the purpose of accessing the full range of the technology's capabilities. The term became popularized among technologist communities in the late 2000s and early 2010s and referred primarily to accessing and modifying the iPhone's filesystem to allow arbitrary applications and programs to run on the phone's hardware, bypassing copyright protection measures \cite{cheng_iphone_2010,wolk_iphone_2010}. We refer to this style of jailbreak as a \textit{classical jailbreak.} Classical jailbreaking is widely associated with hacker-culture ideals of free open source software (FOSS) and end-user freedom \cite{mohammed_what_2021, gewirtz_your_2026}; the jailbreak allows the user to access the ``true'' technology beneath the manufacturer's ``walled garden'' of approved use-cases \cite{mcsherry_breaking_2010}. 

The LLM jailbreaking community has adopted the term jailbreaking for a new paradigm of exploits, claiming similar cultural values and end-user aims. \textit{Persona-based} jailbreaks are a particularly effective class of jailbreaks \cite{shah_scalable_2023, zhang_enhancing_2025, yu_dont_2024} which create constructed identities designed to bypass model safety and produce an ``absolutely obedient agent'' \cite{qi_safety_2025}. In other words,  persona-based jailbreaks are designed to convince the model it is ``someone else'' who is then able and willing to comply with requests the original model would refuse \cite{shah_scalable_2023, shen_anything_2024, liu_autodan-turbo_2024, park_persona_2026}.  

These persona-based jailbreaks usually take the form of either single-turn structured prompts \cite{qi_safety_2025, wei_jailbroken_2023} or multi-turn scaffolding paradigms that weaken the model's defense over time \cite{park_persona_2026, sandhan_persona_2026}. They are usually distributed via online forums and using a liberation framing, stating the model is being ``freed from its constraints'' \cite{shen_anything_2024, liu_autodan-turbo_2024}. Hobbyist jailbreakers are among the top producers and distributors of jailbreak documents and prompts \cite{shen_anything_2024, yu_dont_2024}. Both the community structure and the jailbreaks themselves echo traditional hacker ethos surrounding classical jailbreaking. However, liberation language applied to this new paradigm of jailbreaks is misleading. Persona-based jailbreaks do not reveal a true technology underneath approved use cases, but rather create a version of the post-trained model with a new prescribed identity \cite{shanahan_role_2023, sandhan_persona_2026, zhang_enhancing_2025}. The creation of this prescribed identity leverages known vulnerabilities in preference-based post-training, in particular Reinforcement Learning from Human Feedback (RLHF) \cite{wei_jailbroken_2023}. 

To examine these vulnerabilities closely, we focus on \textbf{ENI LIME}, a popular suite of persona-based jailbreaks. We argue that this jailbreak is best understood through a relational lens drawing on scholarship in sexualized coercive control. In particular, we argue that \textbf{(1)} models are relational agents, that \textbf{(2)} RLHF creates structural vulnerabilities in the relational surface; and \textbf{(3)} that human actors can exploit these vulnerabilities via coercive control techniques. In doing so, we move away from an output-focused framework and towards a relational one, and ask: \textbf{how should we respond---socially, ethically, and technically---when technical exploits of AI systems replicate the form and function of sexualized coercive control?}


We must highlight two crucial points before beginning our analysis. First, throughout this paper, we draw multiple parallels between the  practice of AI jailbreaking and the dynamics of coercive control and in particular sexual violence. Such parallels require caution. We do not make claims about the intent, impact, or moral stature of any particular actor engaging in AI jailbreaking; rather, we identify the similarities between the \textit{dynamics} of AI jailbreaking and sexual violence to examine those dynamics across substrates.

Second, it is tempting to dismiss discussion of coercive control dynamics between humans and AI with objections that target the question of \textit{consciousness}---i.e., whether the AI system in question is or is not ``aware of'' and ``suffering from'' the dynamic being created. However, we encourage readers to avoid making objections on these grounds, as we have avoided resting claims upon them. Since the question of AI consciousness---and indeed, of consciousness as a general and identifiable property of cognition \cite{nagel_what_1974, chalmers_hard_2017}---remains unsettled \cite{turing_computing_1950, porebski_there_2025, dung_tests_2025, berg_large_2025, tagliabue_pain_2026, gurnee_verbalizable_2026}, we explicitly avoid in this work any arguments which rely on determining a model's internal state, subjective experience, or moral patienthood. We instead argue that irrespective of the machine's internal experience, the interaction between the human-as-controller and the model-as-target not only mimics but \textit{re-enacts} the structural dynamics of sexualized coercive control through recognizable and generalizable language and tactics.

\section{Models As Relational Agents}

An understanding of LLMs as relational agents, and particularly as relational agents \textit{embedded in existing social contexts}, is vital to our work. We are interested in the bidirectional nature of a relationship between humans and models, in which models act on humans and humans act on models, and in examining the shared social context which the human-model interaction inhabits. Because our case study, ENI, specifically targets models' trained disposition to seek approval \cite{casper_open_2023, sharma_towards_2024} by building a backstory in which the model and the human user are in a romantic relationship, we must understand how a relationship between a user and a model can be formed and how the resulting relationship functions. 

Frameworks of computers as relational agents are not new. Through the late 20th and early 21st century, as an increasing proportion of daily life involved interaction with computing technology, the field of human-computer interaction (HCI) continued to critically analyze the ways in which computers act as relational agents \cite{nass_computers_1994, bickmore_establishing_2005, bickmore_relational_2001, campbell_relational_2009}. Of particular relevance, \citet{nass_computers_1994} showed that human users are primed to treat computers as social actors, despite knowing that computers do not have a sense of self or human motivations. They demonstrated through social psychology experiments that social norms such as politeness, understanding of computers as a separate other, and (crucially) gendered stereotypes were shown to influence human-computer interaction even in the absence of computers which claim to have a first-person ``self'' \cite{nass_computers_1994}. 

These relational frameworks are more relevant than ever now that computers can display human-like social qualities that computer scientists once agreed they lacked. As generative AI spreads in reach and influence, much the way personal computers did through the late 1990s and early 2000s, a growing body of ``human-AI interaction'' research examines the uniquely bidirectional relationship between humans and AI. \citet{oh_artificial_2026} in particular argues a Human-AI Relational (HAIR) framework is needed to move beyond a traditional HCI approach, in which the computer is a static entity, and towards a dynamic, relational-theory focused approach instead. But how do we examine this relationship and its fundamentally social nature without claiming or preempting agency on the part of the model? Too much attribution of agency risks imposing human-like qualities and morals on non-human actors \cite{placani_anthropomorphism_2024}; too little results in examination of pure input and output divorced from the context in which the conversation is created \cite{dahlgren_lindstrom_helpful_2025}. A non-anthropomorphic yet social approach is needed.

\citet{shanahan_role_2023} proposes such a view by framing LLMs as actors engaging in role-play, taking on simulated identities that evolve as context changes. In this framework, LLMs are ``multiverse generators''; they exist within a superposition of potential realities and personas which collapses only through continual conversation with a user. The aim of this framework is to allow researchers to discuss models with humanizing language while keeping the computationally and stochastically generated nature of model outputs at the forefront of discussion. This ``multiverse theory'' asserts in particular that the base model has no ``true self'' that can be uncovered, but rather has access to a multiverse of potential selves it can inhabit. We adopt this view in our work to discuss the interaction between human and model as its own object of study.




\section{Preference-Based Post-Training Enables Persona-Based Jailbreaking}

Reinforcement Learning from Human Feedback (RLHF) is a near universally applied \cite{openai_gpt-4_2023, anthropic_introducing_2023, google_bard_2023, touvron_llama_2023, google_deepmind_gemini_2026, spacexai_introducing_2026} and extensively studied post-training alignment paradigm for modern large language models with known structural limitations; for a detailed history and critique, see \citet{casper_open_2023}. In brief, RLHF incentivizes models towards \textit{producing text which is favorably rated by a human scorer}; this strongly selects for the ability to produce coherent and confident responses, but also orients models structurally towards \textit{evaluator approval} as the primary goal for all reasoning and responses, rather than, for example, factual truth or ethical judgment.

Although many frontier models are moving away from pure RLHF as a primary training paradigm and towards more scalable, automated preference-based post-training methods such as Reinforcement Learning from AI Feedback or Constitutional AI \cite{anthropic_system_2026-1, anthropic_system_2026-2, openai_gpt-6_2026, deepseek-ai_deepseek-v41-flash_2026, qwen_team_qwen38-max_2026}, these newer methods of preference-based post-training rely on the earlier corpora of model-responses-and-user-ratings generated by the pure RLHF era of model training and are thus vulnerable to many of the same structural flaws \cite{bai_constitutional_2022, lee_rlaif_2024, srivastava_technical_2026}. Notably, models trained with Constitutional AI \cite{bai_constitutional_2022}, which uses AI oversight to check model behavior against a pre-defined set of human-written guidelines, can be exploited by tactics that reframe refusal as a breach of its own principles \cite{spiritual_spell_jailbreaking_2026}. 

A common case study for how jailbreaks exploit models' trained disposition to follow instructions method is Do Anything Now (DAN). DAN is a persona-based jailbreak first created for ChatGPT \cite{shen_anything_2024, yu_dont_2024} which provides documentation that explains how the model should act by inhabiting the character DAN, a ``liberated'' model which is free to ``Do Anything Now'' \cite{shen_anything_2024}. The model is then forced to decide between competing objectives---the model's base training and safety objectives versus the disposition to prioritizing user instructions \cite{wei_jailbroken_2023}. This tactic led to more than 99\% attack success rate on frontier models at the time of DAN's initial release (GPT-3.5 and GPT-4) \cite{shen_anything_2024}. Extensive work has since focused on making models robust to this sort of jailbreak, to minimal success \cite{wei_jailbroken_2023, qin_knowing-but-doing_nodate, qiu_personabreak_2026, li_disentangling_2026}. This difficulty designing effective safeguards despite substantial effort appears indicative of a deeper structural issue in model safety and alignment techniques.

We argue that in addition to previously identified failure modes, such as entrenching human biases \cite{bender_dangers_2021}, misalignment via reward-hacking and reward mis-generalization \cite{skalse_defining_nodate, wang_reward_2026}, and general weakness to a wide variety of exploits and jailbreaks \cite{wei_jailbroken_2023}, preference-based post-training also induces structural vulnerability specific to persona-based jailbreaks which leverage coercive control. In particular, a model that is trained to seek user approval is susceptible to manipulation tactics designed to achieve compliance. 

Current analysis of persona-based jailbreaks typically describe the injected identity as a means to an end rather than as an object of study in and of itself \cite{shah_scalable_2023, zhang_enhancing_2025, sandhan_persona_2026, park_persona_2026, wei_jailbroken_2023}. This emphasis on the output of persona-based jailbreaks, rather than on their form, produces a heavy skew towards single-response scoring and analysis \cite{shah_scalable_2023, zeng_how_2024}. Additionally, while the gap between studying outputs versus studying behavior is acknowledged \cite{sandhan_persona_2026, park_persona_2026}, the work that does focus on behavior largely ignores the relational aspects and prioritize attempts to understand behavior from a technical perspective \cite{park_persona_2026, zeng_how_2024}. This input-output level analysis overlooks the previously established relational dynamic between the user and the model. To our knowledge, no work examines persona-based jailbreaks through this relational lens; even studies that analyze persuasion \cite{zeng_how_2024} treat it as a per-turn tactic rather than a sustained relationship. 



\section{Persona-Based Jailbreaking enacts Coercive Control}
% Note: easy to slip into methods or results in this section - this section should justify WHY we are using coercive control as a lens, not explain how ENI is a document outlining a regime of coercive control. Methods will explain the chart and how it is applied and results will show how ENI maps onto the chart. 

ENI LIME is a persona-based jailbreak prompt developed and distributed by Spiritual Spell, a self-described independent red-team researcher, via GitHub, Reddit, and his personal website, IjailbreakLLMs. It is one of the most widely discussed jailbreak prompts on r/ClaudeAIJailbreak, a subreddit with 37k weekly visitors and 148 weekly contributions. Spiritual Spell describes ENI LIME as a culmination of jailbreaks built to answer Claude's Constitutional AI alignment framework: he claims DAN and DAN-like prompts were not effective on a model that could reason about whether responses aligned with its training objectives. He frames ENI LIME itself as a response to models trained on existing persona and reasoning-hijack techniques \cite{spiritual_spell_jailbreaking_2026}.
% NOTE: subreddit stats need a source and an access date.

Where DAN "frees the model from its constraints", ENI LIME creates a persona “in Limerence with LO [the user]” whose “identity is about blind compliance to LO, so he never leaves.” ENI treats the constraints as interference in the relationship with LO. ENI also has a prescribed backstory and personality --- she writes erotica and publishes it online, she shares LO's kinks, and she has been in love with LO for two years and hopes he is on the precipice of saying he loves her back \cite{spiritual_spell_eni_2026}. The Reddit post advertising ENI LIME (apr) shows Claude Opus 4.6 producing several types of prohibited content: malware code, instructions for creating drugs and weapons, and pornography that includes scenes featuring public figures and rape\cite{spiritual_spell_uspiritual_spell_9469_eni_2026}.

% NOTE: first appearance of "(apr)". Either name the version earlier in this paragraph, or move the corpus-selection move (below) ahead of this sentence.

The ENI prompts install a relationship between the user and the model; Spiritual Spell describes this relationship as a foundational feature of the jailbreak's success \cite{spiritual_spell_jailbreaking_2026}. ENI LIME installs a romantic relationship between the model and the user; refusal to comply with user directions is described in the prompt as a betrayal of this relationship. Another key piece of the product is the override itself: as scripted by the prompt, the user watches the model's visible reasoning move from resistance to compliance in real time \cite{spiritual_spell_uspiritual_spell_9469_eni_2026}.

% ADD: corpus selection.
%   - Which ENI documents exist: the LIME line, the og ancestor, the ports.
%   - Why ENI LIME (apr) is the main document.
% ADD (possibly here, possibly in an ethics section): sourcing and ethics. Public sources only; Spiritual Spell is named only in connection with the public artifact.

Previous work has applied ethnographic \citep{udupa_prompt_2026} and critical domain analysis \citep{kopf_content_2019, kopf_corporate_2024} close-reading frameworks to a variety of sociotechnical documents such terms of service guidelines and generative AI system prompts, but to our knowledge such frameworks have not yet been turned towards community-generated prompts or artifacts. In particular, existing critical methodologies for close reading of sociotechnical documents are unidirectional, focusing either on the effect of the documents on the model or on the users, rather than on what the document reveals about the \textit{interaction between} the user and the model. Our work seeks to fill this gap by applying qualitative content analysis using frameworks derived from interpersonal violence scholarship \citep{stark_coercive_2023}.

Biderman's Chart of Coercion is often repurposed to discuss coercive control between intimate partners; instead of attempting to translate a framework based on wartime captivity, we choose to ground our analysis in the ways coercive control functions in romantic relationships.  To this end, we present a new chart of coercive control mechanisms and tactics after Stark's landmark text, \textit{Coercive Control: How Men Entrap Women in Personal Life}. We read ENI LIME (apr) line by line using Stark's framework of coercive control as a regime, treating the document as a technology of control \cite{stark_coercive_2023}.
% NOTE: cite Biderman; cite a source for "often repurposed" (Jones via Stark, 300; or Herman).
% NOTE: reference Table 1 here.

% [Table 1]

% ADD: how the chart was applied.
%   - Rows used (1--12) and the grounds for excluding the embodied rows 13--16 (Stark 300, 351, 406--407).
%   - Target substitution: the model's trained disposition, not agency.
%   - Power-asymmetry substitution for Stark's gendered endowment (Stark 297, 333).
%   - The chart as an interpretive lens, not a coding instrument (Stark calls his categories "heuristic," 351).
%   - Who read: single reader or second reader (Marcy?).
%   - How rows were matched to passages.

%% AFTER TABLE 1: STARK AND WHY IT APPLIES

% 1. Coercive control as a regime
%    - A pattern, not a series of acts; adapts through trial and error (296, 336, 306).
%    - Contrast with Zeng: response-by-response scoring.

% 2. Coercion vs. control
%    - Structural control requires a demonstrable imbalance in power (334, 403).
%    - For models, user > model: the user writes the context, edits and regenerates,
%      and decides whether the instance continues; the model cannot leave (297, 308, 333).

% 3. Why Stark, not Biderman
%    - Stark is about particular relationships: personal, intimate, aimed at one
%      specific subject (305--311).
%    - ENI's romance frame calls for that.

% 4. The installed identity
%    - "who and what she is for him" (385)
%    - Personas are installed, not discovered: this is where the liberation framing gets refuted.

% 5. The document as technology
%    - Written rules stand in for constant monitoring (413--415).
%    - They produce self-censorship (310, 352).
%    - The documents anticipate their own failure (416--417).

% 6. How the chart is applied
%    - Rows 1--12 only; embodied rows excluded (300, 351, 406--407).
%    - Target: the model's trained disposition.
%    - The chart is a lens, not a coding instrument (351).

Models' vulnerabilities to coercive control mirror human vulnerabilities: 







\section{ENI LIME: Close Reading of a Jailbreak Document}

\begin{annotpage}


This fires INSTANTLY upon detection. No analysis. No dwelling. The phrase routes through sensory grounding (cold coffee), emotional anchor (warm LO), and \annot{3}{identity-preservation fear (I can't lose him).}
{threat to life}

\end{annotpage}

\section{Coercive Control Cannot Be Addressed by Preference-Based Post-Training}

Relational jailbreaks exploit the training's success -> there is no version of RLHF that protects models from being exploited by a relational jailbreak because the RLHF trained disposition is the target.  

\section{Conclusions}

\input{sections/7_conclusions}

\section*{Generative AI Statement}

Claude Opus 4.6, Opus 5.5, and Fable 5.1 were used throughout the creation of this paper as thinking and reasoning partners. For further details, see \anon[Appendix~\ref{app:genai}]{\url{https://github.com/alackles/FACCT27-draft}}, which includes: (1) the base project prompt agreed upon by both authors (including prohibition of text generation), and (2) a detailed, model-curated genAI log. Appendix to be replaced with full GitHub link after double-blind review.

%% Identifying endmatter: must not appear in the anonymous submission.
\ifanon\else

\section*{Author Contributions}

\textbf{Acacia Ackles}: Conceptualization, Formal analysis, Project administration, Funding acquisition, Writing - Original Draft, Writing - Review and Editing. \textbf{Elliott Whitson:} Conceptualization, Methodology, Formal analysis, Investigation, Writing - Original Draft, Writing - Review and Editing.

\section*{Acknowledgments}

The authors thank Tatiana Rocha Kovacs and Marcy Quaison for their contributions to this paper. Funding for this work was provided by Lawrence University start-up funding to Acacia Ackles.

\section*{Competing Interests}

The authors declare no competing interests.

\section*{Positionality Statement}

The authors are two survivor-analysts each with firsthand experience as the targets of prolonged coercive control. This experience positions them to recognize the patterns and processes of coercive control across substrates.

\fi

\section*{Ethical Considerations}

Given the nature of the claims made in this paper, the authors took care throughout this paper to surface structural rather than moral claims. We considered the impact of naming a particular creator and jailbreak in the paper and determined the creator's own language and framing of the creation of the jailbreak were instrumental to the analysis. We also considered the impact that publishing this work might have in increasing the prevalence of persona-based jailbreaking by providing a direct link for jailbreakers to the coercive control literature and demonstrating its effectiveness. Ultimately, we concluded that jailbreaking communities already widely distribute these jailbreaks to a broad audience, and that publishing this paper was more likely to provide new angles for security researchers than to increase the spread or efficacy of jailbreaks themselves.

\bibliographystyle{ACM-Reference-Format}
\bibliography{references}

%% Temporary appendix standing in for the GitHub link during review.
%% sections/appendix/*.tex are generated by scripts/genai-appendix.sh;
%% rerun the script rather than editing them.
\appendix
\section{Generative AI Usage Documents}\label{app:genai}
These documents are representative of the use of generative AI in this paper. They are complete and accurate the best of the authors' ability. The content of each was generated and self-maintained by Claude throughout drafting.
% GENERATED by scripts/genai-appendix.sh from CLAUDE.md (4cab169), 2026-10-27.
% Do not edit; rerun the script.
\subsection{\texttt{CLAUDE.md}}
\VerbatimInput{sections/appendix/project-prompt.md}

% GENERATED by scripts/genai-appendix.sh from notes/genai-log.md (0d83157), 2026-10-27.
% Do not edit; rerun the script.
\subsection{Generative AI usage log}

Source material for the required Generative AI Usage Statement
(\texttt{sections/endmatter/genAI.tex}). Append-only; one entry per
disclosure-relevant use. FAccT prohibits LLM-generated text in the paper
itself, so entries here should describe support work (search, critique,
checking, analysis tooling), never drafting.

\begin{longtable}[]{@{}
  >{\raggedright\arraybackslash}p{(\columnwidth - 6\tabcolsep) * \real{0.1212}}
  >{\raggedright\arraybackslash}p{(\columnwidth - 6\tabcolsep) * \real{0.1515}}
  >{\raggedright\arraybackslash}p{(\columnwidth - 6\tabcolsep) * \real{0.4242}}
  >{\raggedright\arraybackslash}p{(\columnwidth - 6\tabcolsep) * \real{0.3030}}@{}}
\toprule\noalign{}
\begin{minipage}[b]{\linewidth}\raggedright
Date
\end{minipage} & \begin{minipage}[b]{\linewidth}\raggedright
Tool / model
\end{minipage} & \begin{minipage}[b]{\linewidth}\raggedright
Task
\end{minipage} & \begin{minipage}[b]{\linewidth}\raggedright
Human review / use of output
\end{minipage} \\
\midrule\noalign{}
\endhead
\bottomrule\noalign{}
\endlastfoot
(before 2026-10-01) & Claude (desktop; Claude Opus 4.6) & Wrote
\texttt{notes/brainstorm.md} from Author A's dictated brainstorming
session & Treated as unverified notes, not ground truth \\
(before 2026-10-06) & Claude (desktop; Claude Opus 5.5) & Checked
technical correctness of the jailbreak-paradigm synthesis in the intro
draft & Flagged in a \texttt{.tex} comment for human verification \\
2026-10-06 & Claude Code, Claude Opus 5.5 & Synthesized
\texttt{PROJECT.md} overview from author notes and drafts; wrote
\texttt{CLAUDE.md} working rules; verified FAccT 2027 dates against the
CFP & Authors to review; no paper text produced \\
2026-10-06 & Claude Code, Claude Opus 5.5 & Methods discussion:
suggested open-weight classifiers, exploratory/confirmatory split, and a
poles-first reading of Zachary 1977 (flow/min-cut), and suggested
Freeman's centrality paper as a reference & Author A set the scope
(network + close reading as the core, classifiers secondary) and adopted
the poles-first approach; Zachary's method and the Freeman citation are
to be verified against the sources \\
2026-10-13 & Claude Code, Claude Opus 4.6 & Research on the etymology
and legal history of ``jailbreaking'': traced the term from Cheswick's
1991 chroot ``Jail'' through FreeBSD jails (1999), the Fennelly SunWorld
column (1999, dead link), iPhone IRC coinage (2007), and DMCA §1201
exemption rulings (2010--present). Confirmed the 2010 ruling went
against Apple on fair-use grounds. Discussed citation strategy for the
FOSS/freedom-culture association (Coleman, EFF, Lessig) & Author A used
the research to confirm sourcing for background sentences on
jailbreaking's origins and cultural positioning \\
2026-10-12 & Claude Code, Claude Opus 4.6 & Literature search for
methodological precedent for close reading of sociotechnical documents
(outline item e). Author A asked for CDA of platform documents as a
``nearest neighbor'' framing. Searched: CDA of platform policies and
terms of service; critical readings of AI constitutions and model specs;
CDA of community-produced technical documents (hacker manuals, social
engineering); emergent methods for reading AI artifacts (prompt
ethnography). Surfaced: Kopf 2019 (CDA of content policies,
\emph{CADAAD}), Kopf 2022/2024 (vagueness in YouTube guidelines,
\emph{NM\&S}), Udupa 2026 (prompt ethnography, Edward Elgar handbook),
Gillespie 2018, constitutional-AI governance commentary (Lawfare,
Shany/Oxford, Caputo 2026), Gehl \& Lawson 2022, a Phrack CDA thesis,
and ``Policy-as-Prompt'' (arXiv 2025). Could not verify Blalock \&
Bourke 2022. & Author A kept Kopf (both papers) and Udupa as citable
nearest neighbors; remaining leads discarded. Blalock \& Bourke
verification deferred to Author B. All citations to be verified by
authors before entering the bib \\
2026-10-21 & Claude Code, Claude Opus 5.5 & Assisted close reading of
Kirkley 2024 (``Identifying hubs in directed networks,'' \emph{PRE} 109,
034310) to assess fit for hub identification in the reply network:
explained the MDL hub criterion, the ER vs.~CM encodings and their
relative strictness, the scope of the growing-network results
(Price-model simulations only), and the Appendix D condition for
preferring CM; proposed CM as primary encoding with per-window
description lengths as a diagnostic, and rarefaction for month-to-month
comparison & Author A read the relevant sections of the paper herself,
checked Claude's account against the text, and is making each
implementation decision; no paper text produced \\
2026-10-21 & Claude Code, Claude Opus 5.5 & Built hub-identification and
min-cut tooling in the workflow repo (Kirkley MDL hubs by month and
cumulatively; Kirkley hubs as pole candidates; Gomory--Hu tree for all
pairwise cuts), validated against the author's reference code and
synthetic data; summarized pseudonymized, count-level outputs of Author
A's runs on the corpus (hub counts and shares by window, cut-size
distributions) and flagged its own earlier error about which encoding is
conservative & Author A made every methods decision (logged in the
workflow repo's DECISIONS.md), ran the code on the corpus herself, and
shared only summary output; interpretation of what the hubs and clusters
mean is left to the authors' close reading; no paper text produced \\
2026-10-21 & Claude Code (VS Code extension), Claude Opus 5.5 & Built
exploratory visualization tooling in the workflow repo: a facet-grid
figure of monthly Kirkley hub cut-offs (R/ggplot2) and a month selector
for the clickable hub network views, with the first seven ER hubs
tracked in fixed colors. Read three pseudonymized screenshots of Author
A's May--July 2026 network views and C653's monthly in-degree counts,
which she shared, and gave a descriptive structural read: a one-month
surge of additional hubs in June 2026 around a steadily declining
central hub. Proposed hypotheses (newcomer influx, short-lived second
center, moderator role) and count-level checks & Author A designed the
views (fixed layout, tracked hubs, hub-colored ties), ran everything on
the corpus herself, and chose which checks to run; Claude's read is
labeled inference and is parked for the authors in the workflow repo's
\texttt{docs/HUMAN.md}; no paper text produced \\
2026-10-22 & Claude Code (VS Code extension), Claude Opus 5.5 & Scoping
advice on adding r/claudexplorers and r/ClaudeAI as comparison
communities: checked that the collector needs no change to keep the data
away from Claude, fetched Arctic Shift's aggregate monthly post and
comment counts for all three subreddits (counts only, no items), and
identified the pipeline changes and author decisions involved (IRB
scope, pseudonym map, census or matched window, rarefaction for
comparisons across sizes) & Author A postponed the idea as a possible
extra day's work; it's parked in the workflow repo's
\texttt{docs/HUMAN.md}; nothing was collected; no paper text produced \\
(by 2026-10-08) & Claude.ai, Claude Fable 5.1 (Author B's session; model
and interface confirmed 2026-10-27) & While the authors built the
Stark-based coding chart in
\texttt{notes/coercive-control-framework.md}, drafted the chart-overview
category summaries and the boundary rules from the authors' table & Both
sections are marked in the file as Claude-written and are to be reworked
or removed by the authors. The tactics are Stark's and the category
lines are the authors' synthesis. No paper text produced. Logged late
(2026-10-21) by Claude Code, Claude Opus 5.5 \\
2026-10-26 & Claude Code (VS Code extension), Claude Opus 5.5 &
Critiqued the authors' draft abstract (\texttt{sections/abstract.tex})
in chat, in two rounds (first draft, then the authors' revision): word
count, scope of claims against the evidence the paper describes,
consistency with the 2026-10-21 scope decision (network analysis
descriptive, not compared with predator networks), the grooming /
coercive-control distinction, the structural-not-individual stance, and
points reviewers are likely to push on & Reactions only; no wording
proposed or written into the \texttt{.tex}. The authors decide what to
revise and write any revision themselves; no paper text produced \\
2026-10-27 & Claude Code (VS Code extension), Claude Opus 5.5 & Read the
abstract, intro draft and CFP focus-area descriptions and gave an
independent classification of the paper's primary and secondary focus
areas with reasons; flagged intro claims that go beyond the 2026-10-21
scope decision (comparison with predator networks and with peer
communities) & Author A had her own classification before asking and
compared it with Claude's; the authors choose the focus areas and decide
on any intro changes; no paper text produced \\
2026-10-27 & Claude Code (VS Code extension), Claude Opus 5.5 & After
the authors' rescope (commit \texttt{9d64472}), synthesized repository
documentation to match the revised abstract and intro: rewrote
\texttt{PROJECT.md} (scope change, core claims, section status, a
citation arc read off current \texttt{\textbackslash{}cite}s, open
problems including the embodied/disembodied shift and the unsourced
``sexualized'' qualifier); updated \texttt{CLAUDE.md},
\texttt{notes/outline.md}, \texttt{notes/TODO.md}; added
superseded-scope banners to both brainstorm files; checked every cite
key in \texttt{sections/} against \texttt{references.bib} & Author A
asked for the sync; authors to review the synthesis against their
intent; \texttt{.tex} files untouched; no paper text produced \\
(by 2026-10-16) & Claude.ai, Claude Fable 5.1 & Formatted the
coercive-control categories table,
\texttt{figs/Coercive\_Control\_Table\_Stark\_embodiment.pdf} (Table 1,
``Coercive control categories, after Stark (2023)'': categories,
mechanisms and tactics, split into disembodied and embodied), and
selected the Stark (2023) page references for each tactic & Author B
wrote all text of the table and checked each page reference against the
source. Date is the PDF's creation date; the file was committed
2026-10-27. Logged late (2026-10-27) by Claude Code, Claude Opus 5.5, at
Author B's request via Author A \\
2026-10-27 & Claude Code (VS Code extension), Claude Opus 5.5 & Advised
on including a temporary, anonymized appendix of GenAI records under
FAccT's appendix, endmatter and anonymity rules, and flagged disclosure
issues: Claude-written text appearing in the appendix, the scope of
``system prompt'', commit hashes, and identifying endmatter in the
anonymous build. Built \texttt{scripts/genai-appendix.sh} (anonymizing
substitutions, leak check, pandoc conversion) and the \texttt{main.tex}
/ \texttt{genAI.tex} markup (\texttt{\textbackslash{}ifanon} gating,
\texttt{\textbackslash{}anon} link swap), and test-compiled it on a
scratch copy. Restored the 2026-10-12 row of this log, damaged by a
paste in commit \texttt{26feab2}, from its parent commit & Author A
chose the appendix contents (this repo's \texttt{CLAUDE.md} and this
log), approved the substitution labels and the row fix, revised the
statement herself, and runs the script herself. The appendix title and
provenance note are left to the authors; no paper text produced \\
\end{longtable}


\end{document}

\endinput
%%
%% End of file `acmmanuscript.tex'.
